\section{Background and Objects}

\begin{frame}{Notation and Core Objects}
    \small
    \begin{itemize}
        \item Input sequence: $x_{1:T}$ and residual activations $h_{\ell}\in\mathbb{R}^{T\times d}$.
        \item Sparse representation at each layer: $s_{\ell}=E_{\ell}(h_{\ell})$, with decoder reconstruction $\hat{h}_{\ell}=D_{\ell}(s_{\ell})$.
        \item Cross-layer mapping: $T_{\ell\to k}(s_{\ell})\approx s_k$.
        \item Logic validators (LTN predicates/rules) output soft satisfaction scores in $[0,1]$.
    \end{itemize}
    \vspace{0.4em}
    \[
        S(Y_t)=\log P_{\mathrm{LM}}(Y_t)+w_{\mathrm{logic}}\cdot \mathrm{Sat}_{\mathrm{LTN}}(Y_t)
    \]
    \footnotesize
    This same object set is reused in G1 (emergence), G2 (steering), and G3 (diagnostics).
\end{frame}

\begin{frame}{Core Concepts Used in This Thesis}
    \small
    \begin{itemize}
        \item \textbf{Distributed representation:} useful concepts are encoded across multiple directions, so feature identity must be tested for stability.
        \item \textbf{Generalization dynamics:} delayed test gains require phase-aware analysis, not only endpoint metrics.
        \item \textbf{Optimization phases:} memorization and later reorganization can co-exist in the same run.
        \item \textbf{Practical implication:} C1/C2/C3 claims require paired behavioral + mechanistic evidence.
    \end{itemize}
\end{frame}

\begin{frame}{Transformer Internals Relevant to Interpretation}
    \small
    \begin{itemize}
        \item Attention acts as content-based routing:
        \[
            \mathrm{Attn}(Q,K,V)=\mathrm{softmax}\!\left(\frac{QK^\top}{\sqrt{d_k}}\right)V
        \]
        \item Residual stream is the main substrate for feature measurement:
        \[
            h_{\ell+1}=h_{\ell}+\Delta^{\text{attn}}_{\ell}(h_{\ell})+\Delta^{\text{mlp}}_{\ell}(h_{\ell})
        \]
        \item SAEs and transcoders operate over this layerwise residual interface.
    \end{itemize}
\end{frame}

\begin{frame}{Decoder Variants and Invariant Measurement Interface}
    \scriptsize
    \begin{table}
        \centering
        \setlength{\tabcolsep}{3pt}
        \begin{tabular}{@{}p{0.22\linewidth}p{0.26\linewidth}p{0.24\linewidth}p{0.23\linewidth}@{}}
            \toprule
            \textbf{Variant} & \textbf{Main change} & \textbf{Invariant object} & \textbf{Transfer risk} \\
            \midrule
            GPT-style dense decoder & Scale, data, recipe & $h_{\ell}$ and $s_{\ell}$ & Feature drift across checkpoints \\
            \addlinespace[0.2em]
            MoE decoder & Expert routing in FFN path & Residual hooks + validators & Route-dependent instability \\
            \addlinespace[0.2em]
            Long-context decoder & Context handling/efficiency & Same layerwise readout & Shifted activation statistics \\
            \addlinespace[0.2em]
            Instruction-aligned decoder & Objective / preference optimization & Same transcoder interface & Behavior shifts without clear internal causality \\
            \bottomrule
        \end{tabular}
        \setlength{\tabcolsep}{6pt}
    \end{table}
\end{frame}

\begin{frame}{Criteria for Actionable Mechanistic Interpretability}
    \small
    \begin{columns}[T,onlytextwidth]
        \column{0.49\textwidth}
        \begin{block}{Faithfulness and Stability}
            \begin{itemize}
                \item Internal measurements must track causal structure.
                \item Features should persist across seeds/checkpoints.
            \end{itemize}
        \end{block}
        \column{0.49\textwidth}
        \begin{block}{Transferability and Actionability}
            \begin{itemize}
                \item Features/validators should remain meaningful across layers.
                \item Explanations must support controllable interventions.
            \end{itemize}
        \end{block}
    \end{columns}
    \vspace{0.3em}
    \footnotesize
    In this proposal, these criteria are operationalized with reconstruction, sparsity, alignment, validator consistency, and steering-locality metrics.
\end{frame}

% \begin{frame}{Visual Concept Map}
%     \conceptplaceholder
%         {Unified architecture overview}
%         {Diagram to be added: token flow through model layers, SAE hooks per layer, cross-layer transcoder edges, and LTN validator heads feeding both analysis metrics and inference-time steering signals.}
%     \vspace{0.4em}
%     \footnotesize Placeholder created because this specific integrated diagram does not yet exist in \texttt{infufrgs/figures}.
% \end{frame}
